\begin{proof}[Proof of \cref{obj:thm:observable-upper}]
\begin{enumerate}
\item
The assignment--audit product design is a probability measure because its factors are products of Bernoulli probability laws with parameters \(1/2\) and \(q\in[0,1]\). For fixed \(n,d,q\), the branch condition in \cref{obj:def:upper-rule} is constant on the record space. On its positive-retention branch, the score is a finite sum of outcome coordinates multiplied by functions of the finite graph and assignment coordinates; clipping is continuous. The other branch is constant. Thus \(T^{\mathrm{up}}\) is real-valued and measurable on the complete record space.

Its deterministic extension to the randomized-estimator domain is
\[
(O,\xi)\longmapsto T^{\mathrm{up}}(O).
\]
This is an admissible estimator in \cref{obj:def:risk}, so
\[
R_{n,d}(q)
\le
\sup_{\theta\in\mathcal M_{n,d}}
\mathcal L(T^{\mathrm{up}};\theta,q).
\]
% lean: CausalSmith.Experimentation.ThinnedgraphAdditiveRiskFrontier.thinnedDesign_probabilityDesign
% lean: CausalSmith.Experimentation.ThinnedgraphAdditiveRiskFrontier.upperRule_measurable
% lean: Causalean.Stat.minimaxValueENNReal_le_worstCaseRisk

\item
Fix a schedule
\[
\theta=(G,a,t,b)\in\mathcal M_{n,d}.
\]
Because the estimator ignores the independent uniform seed, integration over that seed gives
\[
\begin{aligned}
\mathcal L(T^{\mathrm{up}};\theta,q)
&=
\mathbb E_{P_{\theta,q}}
\left[
\int_0^1
\bigl(T^{\mathrm{up}}(O)-\tau(\theta)\bigr)^2\,d\xi
\right]\\
&=
\mathbb E_{P_{\theta,q}}
\left[
\bigl(T^{\mathrm{up}}(O)-\tau(\theta)\bigr)^2
\right].
\end{aligned}
\]
We bound this loss separately on the two branches of \cref{obj:def:upper-rule}.
% lean: CausalSmith.Experimentation.ThinnedgraphAdditiveRiskFrontier.sqLoss_seedless

\item
First suppose
\[
q>0
\qquad\text{and}\qquad
u(n,d,q)<1.
\]
The product design satisfies \cref{obj:ass:assignment,obj:ass:audit,obj:ass:design-independence}. Together with the population, degree, and schedule hypotheses, these conditions permit application of \cref{obj:lem:score-audit}. It supplies
\[
\mathbb E_{P_{\theta,q}}[T_q]=\tau(\theta)
\]
and hence
\[
\begin{aligned}
\mathbb E_{P_{\theta,q}}
\left[(T_q-\tau(\theta))^2\right]
&=\operatorname{Var}_{P_{\theta,q}}(T_q)\\
&=\operatorname{Var}_{Z}(T_G)
+\frac{4(1-q)}{n^2q}
\sum_{i\in V}
\bigl|\{j:(j,i)\in G\}\bigr|
\,\mathbb E_Z[Y_i(Z)^2],
\end{aligned}
\]
as well as
\[
\operatorname{Var}_{Z}(T_G)\le\frac{5(d+1)^2}{n}.
\]

For each \(i\in V\) and binary assignment \(z\in\{0,1\}^V\), the additive outcome formula and coefficient-mass bound give
\[
\begin{aligned}
|Y_i(z)|
&=
\left|a_i+t_i z_i+\sum_{j:(j,i)\in G}b_{ij}z_j\right|\\
&\le
|a_i|+|t_i z_i|
+\sum_{j:(j,i)\in G}|b_{ij}z_j|\\
&\le
|a_i|+|t_i|
+\sum_{j:(j,i)\in G}|b_{ij}|\\
&\le1,
\end{aligned}
\]
where the penultimate inequality uses \(z_j\in\{0,1\}\), and the last uses \cref{obj:ass:coefficient-mass}. Consequently,
\[
\mathbb E_Z[Y_i(Z)^2]\le\mathbb E_Z[1]=1.
\]
The in-degree bound therefore yields
\[
\begin{aligned}
\sum_{i\in V}
\bigl|\{j:(j,i)\in G\}\bigr|
\,\mathbb E_Z[Y_i(Z)^2]
&\le
\sum_{i\in V}\bigl|\{j:(j,i)\in G\}\bigr|\\
&\le\sum_{i\in V}d
=nd.
\end{aligned}
\]
Since \(n>0\), \(q>0\), and \(q\le1\),
\[
\frac{4(1-q)}{n^2q}\ge0.
\]
Combining the preceding bounds gives
\[
\begin{aligned}
\mathbb E_{P_{\theta,q}}
\left[(T_q-\tau(\theta))^2\right]
&\le
\frac{5(d+1)^2}{n}
+\frac{4(1-q)}{n^2q}\,nd\\
&=
\frac{5(d+1)^2}{n}
+\frac{4d(1-q)}{nq}
=u(n,d,q).
\end{aligned}
\]
All expectations here are integrable because the assignment--audit space is finite.
% lean: CausalSmith.Experimentation.ThinnedgraphAdditiveRiskFrontier.abs_potentialOutcome_le_one
% lean: CausalSmith.Experimentation.ThinnedgraphAdditiveRiskFrontier.auditScore_squared_error_le

\item
To control clipping, first bound the target. Define the constant assignments by
\[
\mathbf1=(1)_{i\in V},
\qquad
\mathbf0=(0)_{i\in V}.
\]
For every \(i\in V\),
\[
Y_i(\mathbf1)-Y_i(\mathbf0)
=t_i+\sum_{j:(j,i)\in G}b_{ij},
\]
so
\[
\begin{aligned}
|Y_i(\mathbf1)-Y_i(\mathbf0)|
&\le |t_i|+
\left|\sum_{j:(j,i)\in G}b_{ij}\right|\\
&\le |t_i|+\sum_{j:(j,i)\in G}|b_{ij}|\\
&\le1.
\end{aligned}
\]
The last inequality follows from \cref{obj:ass:coefficient-mass} and \(|a_i|\ge0\). Averaging these contrasts gives
\[
\begin{aligned}
|\tau(\theta)|
&=
\frac1n
\left|\sum_{i\in V}
\bigl(Y_i(\mathbf1)-Y_i(\mathbf0)\bigr)\right|\\
&\le
\frac1n\sum_{i\in V}
|Y_i(\mathbf1)-Y_i(\mathbf0)|\\
&\le\frac1n\sum_{i\in V}1
=1.
\end{aligned}
\]
Here \(n\ge4\) ensures \(n>0\).
% lean: CausalSmith.Experimentation.ThinnedgraphAdditiveRiskFrontier.abs_tte_le_one

\item
For arbitrary real values with domains
\[
v_{\mathrm{clip}}\in\mathbb R,
\qquad
s_{\mathrm{target}}\in[-1,1],
\]
clipping satisfies
\[
\left|\max\{-1,\min\{1,v_{\mathrm{clip}}\}\}
-s_{\mathrm{target}}\right|
=
\begin{cases}
s_{\mathrm{target}}+1,
&v_{\mathrm{clip}}<-1,\\
1-s_{\mathrm{target}},
&v_{\mathrm{clip}}>1,\\
|v_{\mathrm{clip}}-s_{\mathrm{target}}|,
&-1\le v_{\mathrm{clip}}\le1.
\end{cases}
\]
In the first case,
\[
s_{\mathrm{target}}+1
\le s_{\mathrm{target}}-v_{\mathrm{clip}}
=|v_{\mathrm{clip}}-s_{\mathrm{target}}|,
\]
and in the second,
\[
1-s_{\mathrm{target}}
\le v_{\mathrm{clip}}-s_{\mathrm{target}}
=|v_{\mathrm{clip}}-s_{\mathrm{target}}|.
\]
The third case gives equality. Squaring therefore proves
\[
\bigl(\max\{-1,\min\{1,v_{\mathrm{clip}}\}\}
-s_{\mathrm{target}}\bigr)^2
\le
(v_{\mathrm{clip}}-s_{\mathrm{target}})^2.
\]
Apply this inequality pointwise with \(v_{\mathrm{clip}}=T_q(O)\) and \(s_{\mathrm{target}}=\tau(\theta)\). By the selected branch of \cref{obj:def:upper-rule},
\[
\begin{aligned}
\mathcal L(T^{\mathrm{up}};\theta,q)
&\le
\mathbb E_{P_{\theta,q}}
\left[(T_q-\tau(\theta))^2\right]\\
&\le u(n,d,q)
=\min\{1,u(n,d,q)\}.
\end{aligned}
\]
The nonnegative real expectation agrees with the extended nonnegative loss used in the risk definition.
% lean: CausalSmith.Experimentation.ThinnedgraphAdditiveRiskFrontier.clipped_sq_error_le
% lean: CausalSmith.Experimentation.ThinnedgraphAdditiveRiskFrontier.observable_upper

\item
Now suppose the branch condition fails:
\[
\neg\bigl(q>0\ \text{and}\ u(n,d,q)<1\bigr).
\]
If \(q>0\), this implies \(u(n,d,q)\ge1\). If \(q\) is not positive, the hypothesis \(q\in[0,1]\) gives \(q=0\), and \cref{obj:def:upper-envelope} gives
\[
u(n,d,0)=+\infty\ge1.
\]
Thus in either case
\[
\min\{1,u(n,d,q)\}=1,
\qquad
T^{\mathrm{up}}(O)=0.
\]
The target bound derived above then yields
\[
\begin{aligned}
\mathcal L(T^{\mathrm{up}};\theta,q)
&=\mathbb E_{P_{\theta,q}}[\tau(\theta)^2]\\
&\le\mathbb E_{P_{\theta,q}}[1]
=1
=\min\{1,u(n,d,q)\}.
\end{aligned}
\]
The equality \(\mathbb E_{P_{\theta,q}}[1]=1\) uses the probability normalization of the product design.
% lean: CausalSmith.Experimentation.ThinnedgraphAdditiveRiskFrontier.observable_upper

\item
Both branches establish the same bound for every \(\theta\in\mathcal M_{n,d}\). Taking the supremum gives
\[
\sup_{\theta\in\mathcal M_{n,d}}
\mathcal L(T^{\mathrm{up}};\theta,q)
\le\min\{1,u(n,d,q)\}.
\]
Together with the admissibility inequality established at the outset, this proves
\[
R_{n,d}(q)
\le
\sup_{\theta\in\mathcal M_{n,d}}
\mathcal L(T^{\mathrm{up}};\theta,q)
\le
\min\{1,u(n,d,q)\}.
\]
% lean: Causalean.Stat.worstCaseRiskENNReal_le
\end{enumerate}
\end{proof}
